\clearpage
\section{Теоретические основы и обзор предметной области}
\subsection{Системы управления персональными знаниями}

В эру информационной перегрузки
необходимо эффективно управлять данными.
Традиционные методы хранения и обработки информации,
такие как рукописные заметки или их цифровые аналоги,
хороши для повседневного использования, но
не способны обеспечить необходимую функциональность
для специалистов, работающих с огромными массивами информации.
Возрастающая актуальность этой проблемы
стимулировала развитие специальных систем для управления персональными знаниями
(Personal Knowledge Management, PKM).


В основу PKM-систем легли результаты множества исследований
из области информационной архитектуры и когнитивной психологии.
В последнее время большое влияние оказывает
концепция «Второго мозга» (Second Brain),
популяризированная исследователем Тиаго Форте~\cite{Forte_2022}.
Эта концепция позволяет переложить хранение и организацию информации на внешние носители,
тем самым разгружая рабочую память человека.

Человеческое мышление ассоциативно,
в нем одна концепция может принадлежать разным областям знания.
Иерархические структуры не обладают достаточной гибкостью,
так как требуют однозначной категоризации,
что противоречит устройству человеческого мышления.
Их все чаще заменяют на
сетевые структуры, наподобие
Zettelkasten («картотечный ящик»)~\cite{Ahrens_2017}.
Этот подход, разработанный социологом Никласом Луманом,
заключается в хранении небольших атомарных заметок,
связанных друг с другом.
Связи в нем представляют основную ценность:
информация, хранящаяся в заметках, не абстрагирована,
а встроена в общий пласт знаний,
что ускоряет поиск и позволяет делать сложные,
на первый взгляд, неочевидные умозаключения.

Популярным представителем современной реализации подхода Zettelkasten
является программа Obsidian~\cite{Obsidian}.
В ней заметки хранятся в виде файлов формата Markdown,
которые пользователь может разместить в удобном ему месте.
Заметки поддерживают двунаправленное связывание путем
упоминания заметки в тексте (формат [[wikilink]]).
Кроме того, в заметках можно добавлять специальные пометки, теги,
по которым можно осуществлять быструю фильтрацию и поиск.
Такая архитектура облегчает организацию данных,
а также обеспечивает их суверенитет, поскольку
пользователь сам выбирает место их хранения.

В основе сетевых PKM-систем лежит концепция персонального графа знаний.
Граф знаний представляет собой структуру,
объединяющую информацию о сущностях и
взаимосвязи между ними.
Его узлы~--- это сущности, например, объекты, люди, термины, идеи и так далее.
Ребра~--- направленные типизированные связи.

\begin{figure}[H]
    \centering
    % Масштабируем рисунок точно по ширине текста
    \resizebox{\textwidth}{!}{%
        \begin{tikzpicture}[
                entity/.style={circle, draw, thick, minimum size=2.5cm, align=center, font=\bfseries},
                literal/.style={rectangle, rounded corners, draw, thick, minimum height=1cm, align=center},
                edge/.style={->, >=Stealth, thick},
                label/.style={fill=white, font=\scriptsize, inner sep=1.5pt}
            ]

            % --- УЗЛЫ (Nodes) ---
            % Синий кластер
            \node[entity, fill=blue!10, draw=blue!80] (bob) {Боб};
            \node[entity, draw=blue!80, above left=2.5cm and 2.5cm of bob] (alice) {Алиса};
            \node[literal, draw=blue!80, below left=3cm and 1.5cm of bob] (person) {Человек};
            \node[literal, draw=blue!80, below right=3cm and 1.5cm of bob] (date) {14 июля 1990};

            % Красный кластер (чуть-чуть сдвинули обратно, чтобы сэкономить ширину)
            \node[entity, fill=red!10, draw=red!80, right=4cm of bob] (monalisa) {Мона Лиза};

            % Зеленый кластер
            \node[entity, draw=green!60!black, above right=2.5cm and 1.5cm of monalisa] (davinci) {Леонардо \\ да Винчи};
            \node[literal, draw=green!60!black, below right=2.5cm and 1.5cm of monalisa] (video) {Видео о \\ Джоконде};

            % --- РЕБРА (Edges) ---
            \draw[edge, blue!80] (bob) -- node[label, sloped] {друг} (alice);
            \draw[edge, blue!80] (bob) -- node[label, sloped] {интересуется} (monalisa);
            \draw[edge, blue!80] (bob) -- node[label, sloped] {экземпляр класса} (person);
            \draw[edge, blue!80] (bob) -- node[label, sloped] {дата рождения} (date);

            \draw[edge, green!60!black] (monalisa) -- node[label, sloped] {автор} (davinci);
            \draw[edge, green!60!black] (video) -- node[label, sloped] {предмет} (monalisa);

        \end{tikzpicture}%
    }
    \caption{Пример графа знаний}
    \label{fig:KG_example}
\end{figure}

Подобная структура данных позволяет создавать
кластеры (сообщества) из семантически близких сущностей,
осуществлять поиск на основе связей.
Это имеет особую ценность в системах управления знаниями, так как
превращает разрозненные пользовательские заметки в
связные группы, что делает поиск и работу с данными эффективнее и результативнее.

\subsubsection*{Простые и семантически типизированные связи}

Как правило, в ГЗ все ребра имеют направление и тип,
который отражает характер отношения между сущностями.
Это обеспечивает более продвинутую навигацию, позволяя
отсеивать нерелевантные направления на раннем этапе,
что открывает возможность использовать многошаговые запросы.

В механизме wikilink-ссылок, который использует Obsidian,
по умолчанию нет поддержки типизации.
Однако она может быть добавлена с помощью популярного расширения
DataView~\cite{dataview_plugin}, которое позволяет задавать тип связи при помощи следующего синтаксиса:
\texttt{link\_type:: [[wikilink]]}.

\subsubsection*{Специфика целевой аудитории}

Основной аудиторией систем управления знаниями
являются люди, чья деятельность
тесно сопряжена с обработкой больших массивов данных.
Исследователи, журналисты и другие специалисты
используют PKM системы для агрегации данных, создавая краткие выжимки,
для их трансформации, находя пробелы в знаниях или данных,
и множества других целей.
Например, эти системы могут быть использованы студентами
для подготовки ответов на экзаменационные вопросы на основе материалов занятий.

\subsubsection*{Постановка проблемы: Когнитивное «узкое место»}

У ГЗ есть недостаток, заключающийся в ручном процессе извлечения и связывания информации.
На начальных этапах создание ссылок и их типизация
происходит естественным образом и не требует особых усилий.
Однако при преодолении определенного порога по объему информации,
например, при наличии в хранилище сотен или тысяч заметок
пользователь утрачивает возможность удержать весь контекст в голове.
Из-за этого процесс связывания становится трудозатратным и неточным:
пользователь начинает пропускать связи, особенно неочевидные.
Как следствие, граф становится разреженным и деградирует качество поиска, так как
из-за отсутствующих связей алгоритму не удается найти необходимую информацию.
В итоге граф начинает терять свою ценность.

Таким образом, алгоритм, способный автоматически
анализировать семантику заметок, находить связи между ними
и типизировать их, тем самым поддерживая актуальность и плотность
графа знаний, устранит это «узкое место»
и освободит когнитивные ресурсы пользователей под другие задачи.
